% s-352-p352-m002; original PDF 352, 353, 354
\Aphorism{s-352-p352-m002}{404}
Wenn man das Schweben oder Wirbeln der Staubtheilchen im Sonnenschein betrachtet, dieses nach geheimen Gesetzen bewegte Gestöber feiner Lichtpunkte, so bedarf es keiner allzubeflügelten metaphorischen Phantasie, hier einen Sternenhimmel \Foreign{en miniature} vor sich zu sehen und jedes dieser lichtreflectirenden Staubatome, als Planeten des verkleinerten Sonnensystems, mit mikroskopisch winzigsten Bewohnern zu bevölkern. Diese Raumabstraction, der Schluss vom ausserordentlich Grossen aufs ausserordentlich Kleine und die Erkenntniss der Relativität alles Raummasses, wird jeder nicht ganz bleischweren und pechzähen Intelligenz gelingen; man muss aber einen Schritt weiter dichten und auch für die \Sperr{Zeit} solche Möglichkeiten offen lassen. Es hätte sogar etwas für sich, jenen hypothetischen Wesen, die ein Staubmolecül als ihre Mutter Erde bewohnen, eine entsprechend verkleinerte Zeit anzuweisen, derart dass in einer Secunde unserer menschlichen Zeit Myriaden Generationen solcher mikroskopischer Geschöpfe entstehen und absterben können. Damit wäre ausgeschlossen, dass wir, selbst mit den allerschärfsten Mikroskopen, von den Vorgängen dieser Miniaturwelt je etwas zu sehen bekämen; ein solches Gewimmel und Geflimmer, unvergleichlich geschwinder als die Entladungsvibrationen einer Leydener Flasche oder als die Längsschwingungen eines elastischen Stabes, würde uns als Ruhe erscheinen — sowie etwa umgekehrt ein Riese, dessen vergröberte Elementarempfindung ganze Jahre menschlicher Zeitrechnung dauerte, statt unserer ergrünenden, abwelkenden, erstarrenden und wieder ergrünenden Erde etwas gleichmässig Graues und Unbelebtes erblicken würde. Nach der Atomphysik ist das, was wir einen festen Körper nennen, auch nur ein unendlich rasch periodisch bewegtes Molecularsystem, ein Sternenhaufen, in dem Entfernungen und Umlaufszeiten jedes menschliche Mass \Sperr{unter}schreiten; ja, es ist eine erkenntnisstheoretische Frage, ob wir absolut Ruhendes überhaupt wahrnehmen würden, und empirische Ruhe ist wahrscheinlich nur eine andere Art, nicht eine andere Quantität der Bewegung. Wenn ein Bleigewicht zur Erde fällt, so verliert es seine Geschwindigkeit und bleibt liegen; diese verlorene Geschwindigkeit, nehmen wir an, geht in Molecularbewegung über und erscheint uns als Wärme, aber den minimalen Bewohnern eines betheiligten Molecüls könnte sie als kreisendes Planetensystem und Harmonie der Sphären erscheinen! Es schwindelt Einen zu denken, was für Unendlichkeiten von Raum und Zeit hier unserer Wahrnehmung durch die Maschen schlüpfen: welche Licht- und Tonwelten, welche Vegetation, welche Organismen, welche Historien und Culturen ein einziges jener Sonnenstäubchen während seines Vorübergleitens vor unserem Auge beherbergen mag! Und umgekehrt, was mögen unsere Monde, Erden, Sonnen und Milchstrassen in der Weltanschauung eines gigantischen Wesens bedeuten, dessen Blutkörperchen Billionen Cubikmeilen füllen und dessen Nervenvibrationen Jahrtausende zählen? Vielleicht auch nur ein durch Stoss erwärmtes Bleigewicht, oder ein Stück Fensterscheibe, oder ein Fingerglied, oder ein in der Sonne tanzendes Staubfäserchen? — Genug: begreifen wir nur, wieviele Bewusstseinswelten neben einander hergehen und, durch die blosse Relativität von Zeit- und Raummass, aus dem Einen Stoffe geformt werden können, der uns als unsere Welt, dem Elephanten als seine Welt und der Mücke als ihre Welt erscheint.
\par

